duit aussitôt que les éléments de $\mathfrak m$ sont non inversibles dans $G$. D'autre part, si $(x,y)\notin\mathfrak m$, et si par exemple $x\notin\mathfrak r$, $f(x)$ n'appartient pas à l'idéal maximal de $A$, donc, puisque $g(y)=f(x)$, et que $g$ est local, on a aussi $y\notin\mathfrak s$; on en
\oldpage[IV-4]{195}
conclut que $x$ et $y$ sont inversibles, donc il en est de même de $(x,y)$, ce qui achève de prouver (ii).
  \item[(iii)] L'image de $\mathfrak m$ par l'homomorphisme surjectif $G\to E$ est l'idéal maximal $\mathfrak r$ de $E$, donc l'image de $\mathfrak m^n\cap\mathfrak J'$ est $\mathfrak r^n\cap\mathfrak J$, et comme $\mathfrak J$ est fermé (donc complet) pour la topologie induite par la topologie $\mathfrak r$-adique de $E$ ($\mathbf{0}_{\mathrm I}$, 7.3.5), $\mathfrak J'$ est complet pour la topologie induite par la topologie $\mathfrak m$-préadique de $G$. D'autre part, $G/\mathfrak J'$, muni de la topologie $\mathfrak m$-préadique, est isomorphe à $F$ muni de la topologie $\mathfrak s$-adique, donc est complet pour la topologie quotient de la topologie $\mathfrak m$-préadique de $G$. On en déduit aussitôt que $G$ est complet pour la topologie $\mathfrak m$-préadique.
\end{enumerate}

Ce lemme étant établi, il résulte du théorème de Cohen (\textbf{0}, 19.8.8) que $B''$ est isomorphe à un quotient d'un anneau local noethérien complet \emph{régulier} $C$. Il résulte alors de (19.1.2) que l'immersion $\Spec(B')\to\Spec(C)$ est régulière; comme par hypothèse il en est de même de l'immersion $\Spec(A)\to\Spec(B')$, on conclut de (19.1.5, (iii)) que l'immersion $\Spec(A)\to\Spec(C)$ est régulière; mais cette immersion s'écrit aussi comme composée $\Spec(A)\to\Spec(B)\to\Spec(C)$. Or comme $B$ est régulier par hypothèse, l'immersion $\Spec(B)\to\Spec(C)$ est régulière (19.1.2); donc il en est de même de $\Spec(A)\to\Spec(B)$ par (19.1.5, (iv)).

\begin{corollary}[19.3.3]
\label{IV.19.3.3}
Soit $X$ un préschéma localement noethérien localement immersible dans un préschéma régulier. Alors l'ensemble des $x\in X$ où $X$ est une intersection complète est ouvert dans $X$ et contient l'ensemble des points réguliers de $X$.
\end{corollary}

On peut en effet se borner au cas où $X=\Spec(B/\mathfrak J)$, où $B$ est un anneau régulier et $\mathfrak J$ un idéal de $B$. L'ensemble des $\mathfrak p\supset\mathfrak J$ dans $\Spec(B)$ tels que $\mathfrak J_{\mathfrak p}$ soit régulier dans $B_{\mathfrak p}$ est alors ouvert (16.9.5), et c'est l'ensemble des points de $\Spec(B)$ où $X$ est une intersection complète en vertu de (19.3.2).

\begin{corollary}[19.3.4]
\label{IV.19.3.4}
Soient $k$ un corps, $X$ un préschéma localement de type fini sur $k$, $k'$ une extension de $k$, $X'=X\otimes_k k'$. Soient $x'$ un point de $X'$, $x$ sa projection dans $X$. Pour que $X$ soit intersection complète en $x$, il faut et il suffit que $X'$ soit intersection complète en $x'$.
\end{corollary}

La question étant locale sur $X$, on peut supposer que $X=\Spec(A)$, où $A$ est une $k$-algèbre de type fini, donc quotient d'une algèbre de polynômes $k[T_1,\ldots,T_n]$, si bien que $X$ est sous-préschéma fermé du schéma régulier $Y=\Spec(k[T_1,\ldots,T_n])$ (\textbf{0}, 17.3.7). Si l'on pose $Y'=Y\otimes_k k'=\Spec(k'[T_1,\ldots,T_n])$, $Y'$ est aussi un schéma régulier. Or, puisque $\Spec(k')$ est fidèlement plat sur $\Spec(k)$, il résulte de (19.1.5, (ii)) et du fait qu'il s'agit de préschémas noethériens que l'immersion $Y\to X$ est régulière au point $x$ si et seulement si l'immersion $Y'\to X'$ l'est au point $x'$. La conclusion résulte donc du critère (19.3.2).

\begin{corollary}[19.3.5]
\label{IV.19.3.5}
Soient $A$, $C$ deux anneaux locaux noethériens, $\rho:A\to C$ un homomorphisme local, $B=C/\mathfrak J$ un anneau quotient de $C$, $k$ le corps résiduel de $A$. On suppose que $\rho$ fait de $C$ un $A$-module plat, et que l'anneau $C\otimes_A k$ est régulier. Alors les conditions suivantes sont équivalentes:
\begin{enumerate}
  \item[\textup{a)}] L'idéal $\mathfrak J$ est transversalement régulier relativement à $A$.
  \item[\textup{b)}] $B$ est un $A$-module plat, et $B\otimes_A k$ est un anneau d'intersection complète.
\end{enumerate}
\end{corollary}

\oldpage[IV-4]{196}

En vertu de (19.2.4), la condition \emph{a)} équivaut à dire que $B$ est un $A$-module plat et que $\mathfrak J\otimes_A k$ (qui s'identifie à un idéal de $C\otimes_A k$ en vertu de la platitude de $B$ sur $A$ ($\mathbf{0}_{\mathrm I}$, 6.1.2)) est régulier dans cet anneau. Comme l'anneau $C\otimes_A k$ est régulier, il revient au même, en vertu de (19.3.2), de dire (lorsque $B$ est un $A$-module plat) que $\mathfrak J\otimes_A k$ est un idéal régulier de $C\otimes_A k$, ou que $B\otimes_A k=(C\otimes_A k)/(\mathfrak J\otimes_A k)$ est un anneau d'intersection complète; d'où le corollaire.

\begin{definition}[19.3.6]
\EGAReaderTarget{ega4.c04.def.19-3-6}{IV.19.3.6}
Soit $f:X\to S$ un morphisme \textup{plat} localement de présentation finie. On dit que $X$ est une \emph{intersection complète relativement à $S$} en un point $x$, si la fibre $f^{-1}(f(x))$ est une intersection complète (absolue) au point $x$. On dit que $X$ est une intersection complète relativement à $S$, ou que le morphisme $f$ est un morphisme d'intersection complète, si $X$ est intersection complète relativement à $S$ en chacun de ses points.
\end{definition}

\begin{proposition}[19.3.7]
\label{IV.19.3.7}
Soient $g:X\to S$, $h:Y\to S$ deux morphismes plats localement de présentation finie, $f:Y\to X$ une $S$-immersion, $y$ un point de $Y$, $x=f(y)$, $s=g(x)=h(y)$. On suppose que la fibre $X_s=g^{-1}(s)$ soit régulière au point $x$. Alors, pour que $f$ soit une immersion transversalement régulière au point $y$, il faut et il suffit que $Y$ soit une intersection complète relativement à $S$ au point $y$.
\end{proposition}

Dire que $f$ est transversalement régulière au point $y$ revient, en vertu de (19.2.4), à dire que $f_s:Y_s\to X_s$ est une immersion régulière au point $y$. Mais comme $X_s$ est régulier au point $x=f(y)$, cela équivaut, en vertu de (19.3.2), à dire que $Y_s$ est intersection complète (absolue) au point $y$, d'où la proposition.

\begin{corollary}[19.3.8]
\label{IV.19.3.8}
Soit $f:X\to S$ un morphisme plat localement de présentation finie. L'ensemble $U$ des $x\in X$ en lesquels $X$ est une intersection complète relativement à $S$ est ouvert dans $X$. Si de plus $f$ est propre, l'ensemble $E$ des $s\in S$ tels que $X_s=f^{-1}(s)$ soit intersection complète en chacun de ses points est une partie ouverte de $S$.
\end{corollary}

Comme $E=Y\setmin f(X\setmin U)$, la seconde assertion résulte de la première et du fait que $f$ est une application fermée. La première assertion est locale sur $X$, donc on peut supposer que $X$ est un sous-préschéma fermé d'un préschéma de la forme $Z=S[T_1,\ldots,T_n]$ ($T_i$ indéterminées). Ce dernier étant lisse sur $S$ (17.3.8), ses fibres $Z_s$ sont régulières, et il revient donc au même, en vertu de (19.3.7), de dire que $x\in U$ ou que l'immersion $X\to Z$ est transversalement régulière au point $x$; on en conclut donc le corollaire par (19.2.2).

\begin{proposition}[19.3.9]
\label{IV.19.3.9}
\begin{enumerate}
  \item[\textup{(i)}] Toute immersion ouverte est un morphisme d'intersection complète.
  \item[\textup{(ii)}] Soit $f:X\to S$ un morphisme plat et localement de présentation finie qui est un morphisme d'intersection complète. Pour tout changement de base $g:S'\to S$, $f'=f_{(S')}:X_{(S')}\to S'$ est un morphisme d'intersection complète. La réciproque est vraie si $g$ est fidèlement plat et quasi-compact.
  \item[\textup{(iii)}] Si $f:X\to Y$ et $g:Y\to Z$ sont deux morphismes plats localement de présentation finie qui sont des morphismes d'intersection complète, il en est de même de $g\circ f:X\to Z$.
\end{enumerate}
\end{proposition}

L'assertion (i) découle aussitôt de la définition (19.3.1). Pour prouver (ii), remarquons d'abord que si $f$ est plat et localement de présentation finie, il en est de
\oldpage[IV-4]{197}
même de $f'$ (2.1.4), la réciproque étant vraie si $g$ est fidèlement plat ((2.5.1) et (2.7.1)); en outre pour tout $s'\in S'$, si $s=g(s')$, il est équivalent de dire que $f^{-1}(s)$ est une intersection complète ou que $f'^{-1}(s')$ est une intersection complète (19.3.4); d'où (ii), compte tenu de ce que $g$ est surjectif lorsqu'il est fidèlement plat.

Enfin, sous les hypothèses de (iii), $g\circ f$ est plat et localement de présentation finie. Par définition (19.3.6), on est donc ramené au cas où $Z$ est le spectre d'un corps, et à prouver alors que les anneaux locaux $\sh{O}_{X,x}$ sont des anneaux d'intersection complète, ce qui est contenu dans le résultat plus général suivant:

\begin{corollary}[19.3.10]
\label{IV.19.3.10}
Soient $f:X\to S$ un morphisme plat localement de présentation finie qui soit un morphisme d'intersection complète, $x$ un point de $X$, $s=f(x)$. Si $\sh{O}_{S,s}$ est un anneau noethérien d'intersection complète, il en est de même de $\sh{O}_{X,x}$.
\end{corollary}

On peut évidemment se borner au cas où $S=\Spec(A)$ avec $A=\sh{O}_{S,s}$. Montrons en outre qu'on peut se borner au cas où l'anneau local $A$ est \emph{complet}. En effet, posons $A'=\widehat A$, $X'=X\otimes_A A'$, et soit $s'$ l'unique point fermé de $S'=\Spec(A')$, qui est l'unique point de $S'$ au-dessus de $s$; si $f'=f_{(S')}:X'\to S'$, la fibre $f'^{-1}(s')$ est isomorphe canoniquement à $f^{-1}(s)$, puisque $A$ et $A'$ ont même corps résiduel (\textbf{I}, 3.6.4); il y a donc un seul point $x'\in X'$ au-dessus de $x$ et de $s'$, et on a par suite $\sh{O}_{X',x'}=\sh{O}_{X,x}\otimes_A A'$. On en déduit que les anneaux locaux $\sh{O}_{X,x}$ et $\sh{O}_{X',x'}$ ont des \emph{complétés isomorphes}, car de façon générale, si $E$ est un anneau local qui est une $A$-algèbre (l'homomorphisme $A\to E$ étant local), le séparé complété $(E\otimes_A\widehat A)^{\widehat{\,}}$ de l'anneau $E\otimes_A\widehat A$ muni de la topologie produit tensoriel est isomorphe au séparé complété de $\widehat E\otimes_{\widehat A}\widehat A\cong\widehat E$, donc au séparé complété $\widehat E$ de $E$ ($\mathbf{0}_{\mathrm I}$, 7.7.1). En vertu de la définition (19.3.1), il revient donc au même de dire que $\sh{O}_{X,x}$ est un anneau d'intersection complète, ou que $\sh{O}_{X',x'}$ est un anneau d'intersection complète.

Supposons donc $A$ \emph{complet}; on peut en outre supposer que $X=\Spec(B)$, $B$ étant quotient d'une algèbre de polynômes $C=A[T_1,\ldots,T_r]$. Comme les anneaux de polynômes sur un corps sont des anneaux réguliers (\textbf{0}, 17.3.7), il résulte de (19.3.7) que l'immersion $X\to Z=\Spec(C)$ peut être supposée transversalement régulière relativement à $S$; donc on peut supposer que $B=C/\mathfrak J$, où $\mathfrak J$ est engendré par une suite régulière $(g_i)_{1\leq i\leq n}$ d'éléments de $C$. D'autre part, puisque $A$ est complet, on peut par hypothèse l'écrire $A'/\mathfrak K$, où $A'$ est un anneau local régulier et $\mathfrak K$ un idéal de $A'$ (\textbf{0}, 19.8.8), et l'hypothèse que $A$ est un anneau d'intersection complète entraîne que $\mathfrak K$ est engendré par une suite régulière $(f_j)_{1\leq j\leq m}$ (19.3.2). Dans l'anneau de polynômes $C'=A'[T_1,\ldots,T_r]$, les éléments $f_j$ ($1\leq j\leq m$) forment encore une suite régulière engendrant l'idéal $\mathfrak K'=\mathfrak K[T_1,\ldots,T_r]$ (\textbf{0}, 15.1.4); comme $C'/\mathfrak K'=C$, on voit que si, pour chaque $i$, $g'_i$ est un élément de $C'$ d'image $g_i$ dans $C$, la suite formée des $f_j$ ($1\leq j\leq m$) et des $g'_i$ ($1\leq i\leq n$) est régulière dans $C'$; si $\mathfrak J'$ est l'idéal de $C'$ qu'elle engendre, on a $B=C'/\mathfrak J'$, d'où la conclusion en vertu de (19.3.2), puisque $C'$ est un anneau régulier (\textbf{0}, 17.3.7).

\oldpage[IV-4]{198}

\subsection{Application: critères de régularité et de lissité pour les préschémas éclatés.}
\EGAReaderTarget{ega4.c04.subsec.19-4}{subsection:IV.19.4}

\begin{env}[19.4.1]
\label{IV.19.4.1}
Soient $X$ un préschéma, $Y$ un sous-préschéma fermé de $X$, défini par un Idéal quasi-cohérent $\mathscr J$ de $\sh{O}_X$. Rappelons (\textbf{II}, 8.1.3) que le $X$-schéma $X'$ obtenu en faisant éclater $Y$ est le préschéma
\[
  X'=\operatorname{Proj}(\mathscr S),\qquad\text{où}\qquad \mathscr S=\bigoplus_{n\geq0}\mathscr J^n.
\]
Si $\mathscr J$ est de type fini, le morphisme structural $f:X'\to X$ est projectif. Sans hypothèse sur $\mathscr J$, le sous-préschéma fermé
\[
  Y'=f^{-1}(Y)
\]
de $X'$ est défini par l'Idéal quasi-cohérent $\mathscr J\sh{O}_{X'}$ de $\sh{O}_{X'}$, canoniquement isomorphe à $\sh{O}_{X'}(1)$; de façon plus précise, on a une suite exacte
\begin{equation}
  0\longrightarrow\sh{O}_{X'}(1)\xrightarrow{\widetilde u_0}\sh{O}_{X'}\longrightarrow\sh{O}_{Y'}\longrightarrow0
  \tag{19.4.1.1}\label{IV.19.4.1.1}
\end{equation}
où $\mathscr J\sh{O}_{X'}$ est l'image de $\widetilde u_0$ (\textbf{II}, 8.1.7 et 8.1.8). Comme au voisinage d'un point de $X'$, le $\sh{O}_{X'}$-Module $\sh{O}_{X'}(1)$ est libre de rang $1$, $\mathscr J\sh{O}_{X'}$ est engendré dans un tel voisinage par une section inversible de $\sh{O}_{X'}$ et $\mathscr J\sh{O}_{X'}/\mathscr J^2\sh{O}_{X'}$ est par suite un $\sh{O}_{X'}/\mathscr J\sh{O}_{X'}$-Module libre de rang $1$. Ceci prouve la première assertion du lemme suivant:
\end{env}

\begin{lemma}[19.4.2]
\label{IV.19.4.2}
L'immersion canonique $Y'\to X'$ est régulière et de codimension $1$, et $Y'$ est canoniquement isomorphe au $Y$-schéma $\operatorname{Proj}(\operatorname{gr}_{\mathscr J}^\bullet(\sh{O}_X))$.
\end{lemma}

La seconde assertion résulte de (\textbf{II}, 3.5.3) appliqué à l'injection canonique $Y\to X$, compte tenu de ce que $\mathscr S\otimes_{\sh{O}_X}\sh{O}_Y=\operatorname{gr}_{\mathscr J}^\bullet(\sh{O}_X)$ par définition, puisque $\sh{O}_Y$ est isomorphe à $\sh{O}_X/\mathscr J$ et $\mathscr J^n\otimes_{\sh{O}_X}(\sh{O}_X/\mathscr J)$ à $\mathscr J^n/\mathscr J^{n+1}$.

\begin{corollary}[19.4.3]
\label{IV.19.4.3}
Si l'immersion canonique $Y\to X$ est quasi-régulière, la restriction $g:Y'\to Y$ du morphisme $f$ est lisse.
\end{corollary}

En effet, $\sh{N}_{Y/X}=\mathscr J/\mathscr J^2$ est alors un $\sh{O}_Y$-Module localement libre de rang fini et $\operatorname{gr}_{\mathscr J}^\bullet(\sh{O}_X)$ est isomorphe à $\mathbf S_{\sh{O}_Y}^\bullet(\sh{N}_{Y/X})$ (16.9.8), donc $Y'$ est $Y$-isomorphe à $\mathbf P(\sh{N}_{Y/X})$ et par suite est lisse sur $Y$ (17.3.9).

Nous verrons plus tard (chap. V) que $g$ est encore lisse lorsque l'immersion $Y\to X$ « ne présente que des singularités ordinaires ».

\begin{proposition}[19.4.4]
\label{IV.19.4.4}
Avec les notations de \textup{(19.4.1)}, supposons que $X$ soit localement noethérien et que le morphisme $g:Y'\to Y$, restriction de $f$, soit lisse. Soient $x'$ un point de $Y'$, $x=g(x')$. Pour que $Y$ soit régulier au point $x$, il faut et il suffit que $Y'$ soit régulier au point $x'$; lorsqu'il en est ainsi, $X'$ est alors régulier au point $x'$.
\end{proposition}

La première assertion résulte de (17.5.8) et la seconde de (19.1.1) et (19.4.2).

\begin{corollary}[19.4.5]
\label{IV.19.4.5}
Sous les hypothèses générales de \textup{(19.4.4)}, supposons que $Y$ soit régulier au point $x$. Alors, pour toute générisation $x_1$ de $x$ dans $X$ n'appartenant pas à $Y$, $X$ est régulier au point $x_1$. Si $\operatorname{Reg}(X)$ est ouvert \textup{((6.12), par exemple si $X$ est \emph{excellent} (7.8.6, (iii)))}
\oldpage[IV-4]{199}
et si $Y$ est régulier, alors il existe un voisinage ouvert $U$ de $Y$ dans $X$ tel que $X$ soit régulier aux points de $U\setmin Y$.
\end{corollary}

Pour prouver la première assertion, on peut se borner au cas où $X$ est un schéma local de point fermé $\boldsymbol{x}$; l'hypothèse entraîne alors que $Y$ est régulier (\textbf{0}, 17.3.2), donc il en est de même de $Y'$ par (19.4.4). Notons maintenant que le morphisme $f:X'\to X$ est projectif, donc propre, et par suite (\textbf{II}, 7.2.1), si l'on considère l'unique point $x'_1$ de $X'$ au-dessus de $x_1$ (\textbf{II}, 8.1.3), il existe une spécialisation de $x'_1$ appartenant à $Y'$; en lui appliquant (19.4.4) et utilisant (\textbf{0}, 17.3.2) on en conclut que $X'$ est régulier au point $x'_1$, donc $X$ régulier au point $x_1$ puisque $X'\setmin Y'$ est isomorphe à $X\setmin Y$ (\textbf{II}, 8.1.3). Si $Y$ est régulier, tout point de $X\setmin Y$ qui est générisation d'un point de $Y$ appartient donc à $\operatorname{Reg}(X)$; si $\operatorname{Reg}(X)$ est ouvert, $U=Y\cup\operatorname{Reg}(X)$ est donc localement constructible et stable par générisation, donc ouvert ($\mathbf{0}_{\mathrm{III}}$, 9.2.5) et répond à la question.

\begin{proposition}[19.4.6]
\label{IV.19.4.6}
Soient $g:X\to S$ un morphisme, $Y$ un sous-préschéma fermé de $X$ défini par un Idéal quasi-cohérent $\mathscr J$, $X'$ le $X$-schéma obtenu en faisant éclater $Y$, $Y'$ l'image réciproque de $Y$ dans $X'$.
\begin{enumerate}
  \item[\textup{(i)}] Les conditions suivantes sont équivalentes:
  \begin{enumerate}
    \item[\textup{a)}] $\operatorname{gr}_{\mathscr J}^\bullet(\sh{O}_X)$ est une $\sh{O}_X$-Algèbre $g$-plate.
    \item[\textup{b)}] Pour tout $n\geq0$, le $\sh{O}_X$-Module $\sh{O}_X/\mathscr J^{n+1}$ est $g$-plat \textup{(en d'autres termes, le \emph{$n$-ème voisinage infinitésimal} $Y^{(n)}$ de $Y$ dans $X$ (16.1.2) est $S$-\emph{plat})}.
    \item[\textup{c)}] Pour tout changement de base $S_1\to S$ et tout entier $n\geq1$, si l'on pose $X_1=X\times_S S_1$, l'homomorphisme canonique $\mathscr J^n\otimes_{\sh{O}_S}\sh{O}_{S_1}\to\mathscr J^n\sh{O}_{X_1}$ est bijectif.
  \end{enumerate}

  Lorsqu'il en est ainsi, $Y'$ est $S$-plat, et pour tout changement de base $S_1\to S$, si $Y_1$ est l'image réciproque de $Y$ dans $X_1$, le préschéma $X'_1$ obtenu en faisant éclater $Y_1$ dans $X_1$ est canoniquement isomorphe à $X'\times_S S_1$.
  \item[\textup{(ii)}] Supposons que les conditions équivalentes de \textup{(i)} soient satisfaites et de plus que les morphismes $X\to S$ et $Y\to S$ soient localement de présentation finie. Alors $\mathscr S=\bigoplus_{n\geq0}\mathscr J^n$ est une $\sh{O}_X$-Algèbre de présentation finie, le morphisme $X'\to X$ est de présentation finie, l'immersion canonique $Y'\to X'$ est transversalement régulière de codimension $1$ relativement à $S$, et l'ensemble des points de $X$ \textup{(resp. $X'$)} où $X$ \textup{(resp. $X'$)} est $S$-plat, est un voisinage de $Y$ \textup{(resp. $Y'$)}.
\end{enumerate}
\end{proposition}

\begin{enumerate}
  \item[(i)] L'équivalence de \emph{a)} et \emph{b)} résulte aussitôt de ($\mathbf{0}_{\mathrm I}$, 6.1.2); pour montrer l'équivalence de \emph{b)} et \emph{c)}, on peut se borner au cas où $S=\Spec(A)$ et $X=\Spec(B)$ sont affines, avec $\mathscr J=\widetilde{\mathfrak J}$, où $\mathfrak J$ est un idéal de $B$; la suite exacte des Tor montre alors que la condition \emph{c)} implique que $\operatorname{Tor}_1^A(B/\mathfrak J^{n+1},A')=0$ pour tout entier $n\geq0$ et toute $A$-algèbre $A'$; en prenant pour $A'$ une algèbre somme directe $A\oplus M$, où $M$ est un $A$-module quelconque et où la multiplication est définie par $(a,m)(a',m')=(aa',am'+a'm)$ on voit que la condition précédente équivaut à $\operatorname{Tor}_1^A(B/\mathfrak J^{n+1},M)=0$, donc au fait que $B/\mathfrak J^{n+1}$ est un $A$-module plat. Si l'on pose $\mathscr J_1=\mathscr J\otimes_{\sh{O}_S}\sh{O}_{S_1}$, on a alors $\bigoplus_{n\geq0}\mathscr J_1^n=(\bigoplus_{n\geq0}\mathscr J^n)\otimes_{\sh{O}_S}\sh{O}_{S_1}$ et l'assertion relative au préschéma éclaté $X'_1$ s'en déduit aussitôt. Enfin, pour prouver que $Y'$ est $S$-plat, on peut encore se borner au cas où $S$ et $X$ sont affines; si l'on pose $C=\operatorname{gr}_{\mathfrak J}^\bullet(B)$, tout point de $Y'$ admet alors, en vertu de (19.4.2) et de (\textbf{II}, 2.3.6), un
\oldpage[IV-4]{200}

voisinage ouvert affine dont l'anneau e